\begin{ExerciseBlock}
\ExerciseHeading{04}{Looping}{work_and_energy_vertical_loop}
Wie in der Abbildung dargestellt, rollt eine Masse $m$ aus einer Höhe $y=h$ auf einer beliebigen Bahnkurve herab und tritt anschließend in einen Looping mit Radius $R=10\,$m ein. Nach Durchlaufen des Loopings rollt die Masse $m$ auf gerader Strecke weiter.

\begin{center}
	\CachedTikzFigure{content/resources/chapters/work_and_energy/tikz/vertical_loop.tex}
\end{center}

\begin{enumerate}
	\item Welche Geschwindigkeit muss die Masse $m$ am höchsten Punkt des Loopings haben, um nicht herabzufallen, und somit den Looping vollständig durchlaufen zu können?

	\item Aus welcher Höhe $h_{min}$ muss die Masse mindestens herunter rollen, um den Looping vollständig zu durchlaufen.

	\item Geben Sie die Winkelgeschwindigkeit $\omega(y)$ der Kreisbewegung im Looping in Abhängigkeit der jeweiligen Höhe $y$.

	\item Welche Geschwindigkeit $v_E$ hat die Masse $m$ nach Durchlaufen des Loopings auf der geraden Strecke?
\end{enumerate}
\end{ExerciseBlock}
\UseExerciseSolution{work_and_energy_vertical_loop}
